\section{Scope, status, and exact models}

This blueprint covers all five conclusions of Theorems 1 and 2 of
\href{https://arxiv.org/abs/2606.25005v2}{arXiv:2606.25005v2}.
It distinguishes three statuses: an ordinary mathematical proof supplied
below, a Lean declaration checked in the small header kernel, and a
geometric construction obligation still requiring a complete proof.
Only nodes corresponding exactly to the checked header kernel carry
\texttt{leanok}. All other Lean names are planned names.
In particular, the universal main theorems are goals until the outstanding
geometric instantiation is established. A certificate theorem is an
intermediate theorem, not a replacement for constructing its certificate.

The new construction changes the boundary control colour in the hierarchical
case. The crown-cover proof uses new positive constants. These changes do
not add convexity, primitivity, recognizability, or lattice assumptions.
We interpret polygonal tiles as compact simple polygonal disks. Collinear
listed vertices are permitted. Rotations are represented by distinct forms
before the construction; only translations act on placed tiles.

\begin{definition}[Polygonal disk]\label{def:polygon}
\lean{Tilings.Geometry.PolygonalDisk}
A polygonal disk is a compact regular-closed set $B\subset\R^2$ with
nonempty interior and a finite simple polygonal boundary. It carries
indexed boundary vertices and nonzero segments. Consecutive collinear
vertices may be retained as data. Its interior is path connected.
\end{definition}

For Lean, use $\operatorname{EuclideanSpace}\R(\operatorname{Fin}2)$.
The usual product metric on $\R\times\R$ is not silently substituted for
the Euclidean metric in the distance and sine estimates. A vertex-only
\texttt{Polygon} interface is not yet a filled polygonal disk; the bridge
between boundary data and the carrier is a separate development.

\begin{definition}[Finite protoset and placements]\label{def:protoset}
\lean{Tilings.Protoset}
\uses{def:polygon}
A protoset is a finite type $P$ with a polygonal disk $B_p$ for each $p\in P$.
A placed tile is $(p,v)$, with carrier $v+B_p$. Colours belong to the form,
so congruent coloured forms remain distinct.
\end{definition}

\begin{definition}[Placed tiles and patches]\label{def:placed}
\lean{Tilings.PlacedTile,Tilings.Patch}
\uses{def:protoset}
Write $|A|$ for the carrier of a placed tile and $|F|$ for the union of the
carriers in a patch. A patch is a finite set of placed tiles with pairwise
disjoint interiors. Translations act on the anchor coordinates.
\end{definition}

\begin{definition}[Plane tilings and side-to-side]\label{def:tiling}
\lean{Tilings.PlaneTiling,Tilings.SideToSide}
\uses{def:placed}
A plane tiling $T$ is a set of placed tiles whose carriers cover $\R^2$ and
whose interiors are pairwise disjoint. It is side-to-side if two tile sides
with positive-length overlap coincide as whole segments. Local finiteness
is a derived property, not an extra hypothesis of the final theorem.

\emph{Plane coverage is part of every theorem statement below.} The
paper's introductory word ``tiling'' (main.tex, line 92) also admits an
arbitrary nonoverlapping set; that usage is kept here only for finite
patches. Properness of a mere nonoverlapping set is not the notion used:
without coverage a single inert decorated tile would be a proper tiling
projecting outside the family, the surjectivity device of
Lemma~\ref{lem:glues} and the reverse direction of the collar conversion
would fail, and nontrivial sides could lack an opposite neighbour.
\end{definition}

\begin{lemma}[Local finiteness and countability]\label{lem:localfinite}
\lean{Tilings.PlaneTiling.locallyFinite,Tilings.PlaneTiling.countable}
\uses{def:tiling}
Every plane tiling by a finite protoset is locally finite and countable.
\end{lemma}
\begin{proof}
Choose an open disk in the interior of each form. The minimum of these
finitely many positive radii is $r>0$, and the maximum of the form diameters
is $D<\infty$. If a tile meets a compact set $K$, its chosen interior disk
centre lies in the closed $D$-neighbourhood of $K$. Centres belonging to
distinct tiles are at distance at least $2r$, since their open interior
disks are disjoint. A compact metric set has only finitely many
$2r$-separated points, by a finite cover with disks of radius $r/2$.
This proves local finiteness.
For countability choose, for each tile, a rational point in its nonempty
open interior. Chosen rational points for different tiles are distinct.
The injection into $\mathbb Q^2$ proves countability. This proof does not
require rational polygon coordinates.
\end{proof}

\begin{definition}[Substitution]\label{def:substitution}
\lean{Tilings.Substitution}
\uses{def:placed}
A substitution $\tau$ has a common scale $\theta>1$ and a finite patch
$\tau p$ with support $\theta B_p$ for each $p\in P$. On a placed tile
$(p,v)$ it produces the translates by $\theta v$ of $\tau p$.
Iterates are defined recursively. A supertile is a translate of $\tau^n p$,
$n\in\N$. A composition is an inverse decomposition of a plane tiling.
Decorated substitutions are allowed to be relations rather than functions.
\end{definition}

\begin{lemma}[Supports and finite parents]\label{lem:finiteparents}
\lean{Tilings.Substitution.iterate_support,Tilings.Substitution.finite_parents}
\uses{def:substitution}
The support of $\tau^n p$ is $\theta^n B_p$. A fixed placed child has only
finitely many possible placed parents under $\tau$.
\end{lemma}
\begin{proof}
The support identity follows by induction, using the exact support identity
for every child and distribution of a homothety over a union.
A possible parent is determined by its form and the child position in that
form's finite substitution patch. If the child's anchor is $v$ and the
chosen child offset is $w$, the parent's anchor is $(v-w)/\theta$.
There are finitely many such pairs of form and position.
\end{proof}

\begin{definition}[Substitution and hierarchical families]\label{def:hierarchy}
\lean{Tilings.IsSubstitutionTiling,Tilings.IsHierarchical,Tilings.IsSideToSideHierarchical}
\uses{def:tiling,def:substitution}
A plane tiling is a substitution tiling if every finite patch in it occurs,
up to translation, in a supertile. It is hierarchical if there are plane
tilings $T_0=T,T_1,\ldots$ with $\tau T_{n+1}=T_n$.
It is side-to-side hierarchical if every $T_n$ is side-to-side.
\end{definition}

\begin{lemma}[Substitution tilings admit substitution compositions]
\label{lem:substitution-hierarchy}
\lean{Tilings.SubstitutionTiling.exists_composition}
\uses{def:hierarchy,lem:finiteparents,lem:localfinite}
Every substitution plane tiling has a composition which is again a
substitution plane tiling.
\end{lemma}
\begin{proof}
Enumerate its tiles $A_1,A_2,\ldots$. Each initial finite set is contained in
a supertile of positive order: enlarge the patch by sufficiently many
tiles if necessary to exclude order zero. Choose the previous-level
supertile as a parent cover. By the finite-parent lemma, thin this sequence
to stabilize the parent of $A_1$, then of $A_2$, and so forth. The diagonal
sequence yields parents $B_i$ such that every finite set of the $B_i$ is
contained in a supertile. Thus they have disjoint interiors and form a
substitution patch at every finite scale. Their decompositions contain
every $A_i$. No extra child can occur: a child has nonempty interior, and
the plane is already covered by $T$ with disjoint interiors. Consequently
the decomposition is exactly $T$; scaling gives coverage of the parent
plane. Iterate this existence statement using dependent choice.
The Lean proof may equivalently use a finite inverse system of partial
parent assignments and K\"onig's lemma.
\end{proof}

\begin{lemma}[Powers and hierarchical tilings]\label{lem:hierarchy-step}
\lean{Tilings.hierarchical_power_iff}
\uses{def:hierarchy}
$\tau$-hierarchy and $\tau^m$-hierarchy coincide for $m>0$.
Side-to-side hierarchy requires a separate check on the inserted levels.
\end{lemma}
\begin{proof}
Take every $m$th level in one direction. In the other, between two levels
insert the $m-1$ intermediate decompositions of the upper tiling. Exact
supports and disjoint interiors ensure that these are plane tilings.
Side-to-side at the two endpoints does not by itself imply side-to-side
at the intermediate levels, which is why the alignment tuples are needed.
\end{proof}

\begin{definition}[Crowns, FLC, and STS]\label{def:crown}
\lean{Tilings.Crown,Tilings.AllowedCrown}
\uses{def:placed,def:substitution,lem:localfinite}
The full pointed crown at a listed tile vertex $v$ consists of all tiles
whose closed carriers contain $v$, with distinguished centre $v$.
A crown is allowed if it occurs in a supertile. Equality up to translation
preserves the distinguished centre.
\end{definition}

\begin{definition}[Finite local complexity]\label{def:flc}
\lean{Tilings.Substitution.FLC}
\uses{def:crown}
FLC means that the set of allowed full pointed crowns, including incomplete
crowns on supertile boundaries, is finite up to translation.
\end{definition}

\begin{definition}[STS hypothesis]\label{def:sts}
\lean{Tilings.Substitution.STS}
\uses{def:tiling,def:substitution}
STS means that every supertile is side-to-side.
\end{definition}


\begin{definition}[Hereditary listed sides]\label{def:hereditary}
\lean{Tilings.Substitution.HereditarySides}
\uses{def:substitution,def:sts}
Every image of a listed prototile vertex is a listed child vertex, and each
scaled parent listed side is a union of complete child listed sides.
This is an extra core hypothesis, not a consequence of STS: an irrational
dummy cut on unit-square sides is not hereditary for a $2\times2$ subdivision.
Finite refinement retaining such a cut cannot repair heredity, since it
would contain the infinite orbit $\{2^n a\}\bmod1$.
The final main theorems remove this extra hypothesis explicitly.
\end{definition}

\begin{lemma}[STS gives FLC]\label{lem:sts-flc}
\lean{Tilings.Substitution.STS.finite_crowns}
\uses{def:sts,def:flc}
Under STS the allowed pointed crowns form a finite set up to translation.
\end{lemma}
\begin{proof}
At a vertex, a tile in the crown is fixed by its form and the listed vertex
identified with the centre. In a face-to-face polygonal patch there is no
vertex strictly inside another tile side with an incident open sector:
the sector on that side would force a partial side overlap. The number of
incident tiles is bounded by the smallest positive angle among the finite
forms. Hence a crown is a subset of a finite list of positioned forms.
Incomplete crowns are subsets of the same list.
The proof of the no-partial-side assertion uses that a supertile fills its
simple polygonal carrier, rather than an arbitrary disconnected patch.
\end{proof}

\begin{definition}[Bounded local rules and finite decorations]\label{def:localrule}
\lean{Tilings.LocalRule,Tilings.Decoration,Tilings.Proper}
\uses{def:tiling}
A local rule is a possibly infinite set $\mathcal R$ of finite forbidden
patches, together with $D<\infty$ bounding all their support diameters.
A full plane tiling is proper if it contains no translated occurrence of
a member of $\mathcal R$. A decoration is a finite decorated form type,
with a surjection to the original forms preserving carriers after anchor
normalization. Its induced projection acts tilewise.
\end{definition}

\begin{definition}[Soficity]\label{def:sofic}
\lean{Tilings.IsSofic,Tilings.IsSideToSideSofic}
\uses{def:localrule}
A family is sofic if it is exactly the projections of the proper plane
tilings for a finite decoration and a bounded local rule.
It is side-to-side sofic if the rule is equality of colours on complete
matching sides, together with side-to-side adjacency.
\end{definition}

This is the continuous geometric notion in the paper. It is not replaced
by a $\mathbb Z^2$ subshift of finite type. The forbidden set may be infinite;
the decoration alphabet and the radius are finite. Both quantifiers in the
definition of soficity range over plane tilings only, as fixed in
Definition~\ref{def:tiling}.

\begin{lemma}[Finite endofunction has an idempotent power]
\label{lem:finite-idempotent}
\lean{Tilings.FiniteMap.exists_idempotent_iterate}
For a function $c$ on a finite set there is $\ell>0$ with $c^{2\ell}=c^\ell$.
\end{lemma}
\begin{proof}
Finitely many functions imply $c^a=c^{a+b}$ for some $a\geq0,b>0$.
Then $c^{a+i}=c^{a+b+i}$ for every $i\geq0$. Choose a positive multiple
$\ell$ of $b$ with $\ell\geq a$. Both $\ell$ and $2\ell$ lie in the periodic
tail and have the same residue modulo $b$, giving the equality.
Restricting the image of $c^\ell$ makes every retained image point fixed
by $c^\ell$. This is the central-form stabilization used below.
\end{proof}
